
\documentstyle[12pt]{article}

\title{OSCILLATORS IN RESONANCE}
\author{
 Antonio Elipe ${}^1$ and
Andr\'e Deprit${}^{1,2}$  \\
{\small
${}^{1}$Grupo de Mec\'anica Espacial. Universidad de Zaragoza,
50009 Zaragoza. Spain}\\
{\small
${}^{2}$National Institut of Standards and Technology,
Washington D.C., U.S.A. }
 }
 \date{}
\pagestyle{empty}
\begin{document}
\maketitle

In this communication, we  consider dynamical systems defined by
Hamiltonians of the type
\[
{\cal H} = {\cal H}_0 + {\cal H}_1;
\] the principal part consists of a finite number of harmonic
oscillators
\begin{equation}\label{hamCar} {\cal H}_0 = \frac{1}{2}{\displaystyle
\sum _{1\le j \le n}} (X_j^2 + \omega_j^2 x_j^2)
\end{equation} while ${\cal H}_1 ({\mbox{$\bf x$}}, {\mbox{$\bf X$}})$
is a weak perturbation depending on the coordinates
${\mbox{$\bf x$}}= (x_1,\ldots,x_n)$ and of the moments ${\mbox{$\bf
X$}}= (X_1,\ldots,X_n)$.
\par A finite family $(1\le j\le n)$ of canonical transformations
\begin{equation}\label{Poincare}
\begin{array}{l} x_j = [(2 \, \Phi_j)/(\omega _j)]^{1/2} \sin
\varphi_j , \\[1.5ex]
	X_j = (2 \, \omega _j\, \Phi_j)^{1/2} \cos \varphi_j
\end{array}
\end{equation} due to Poincar\'e converts  (\ref{hamCar}) into the
function
\begin{equation}\label{hamPoin} {\cal H}_0 =  {\displaystyle
\sum_{1\le j \le n} }\omega_j\, \Phi_j .
\end{equation} The perturbation is now a function of the angles
${\mbox{$\bf \varphi$}}= (\varphi_1,\ldots,\varphi_n)$ and the actions
${\mbox{$\bf \Phi$}}= (\Phi_1,\ldots,\Phi_n)$.  To a General Theory of
Perturbation, Poincar\'e's variables create the usual difficulty: the
perturbation may introduce zero divisors when the frequencies
$(\omega_i)$ are commensurable.
\par For pairs of harmonic oscillators, except for the resonances 1:1
(Kummer \cite{Kummer}) where the difficulty has been solved exactly
\cite{Dep91}, authors follow the example of Contopoulos
\cite{ContMout} and Verhulst \cite{Verhulst} by assuming that the
ration $\omega_1/\omega_2$ differs from the resonance by an
infinitesimal detuning. This stratagem allows them to use the Birkhoff
normalization or an equivalent averaging method near the resonance.
\par In this communication, we propose a change of variables that will
handle the commensurabilities without recourse to a detuning.
\par It is based on the Lissajous variables introduced in \cite{Dep91}
for a pair of oscillators of equal frequencies. The latter proved
useful in Galactic Dynamics  and more
generally in Molecular Physics; they were
even arranged in \cite{Ferrer} to accommodate some situations in
resonances (1:1:1) for three oscillators. But the resonance (1:1) is
an exceptional event, and Molecular Physics is interested in a general
theory involving any number of oscillators in any type of resonances.
To give one example, we mention the famous Fermi resonance (2:1).

\begin{thebibliography}{}

\bibitem{Kummer}
M. Kummer, in {\sl Local and global methods in nonlinear dynamics,} edited
by A.V. S\'aenz et
al. LNP 252 (Springer, Silver Spring, 1979), p. 19.

\bibitem{ContMout}
	G. Contopoulos and M. Moutsoulas,
	Astron.~J.  {\bf 70}, 817  (1965).

\bibitem{Verhulst}
F. Verhulst, Philos. Trans. R. Soc. A. {\bf 290}, 435 (1979).

\bibitem{Dep91}
  A. Deprit, Celest. Mech. {\bf 51}, 201 (1991).


\bibitem{Ferrer}
S. Ferrer and J. G\'arate, in  {\sl New Trends for Hamiltonian Systems and
Celestial
Mechanics}, edited by E. A. Lacomba and J. Llibre, (World Scientific,
Singapore, 1996) p.
179.

\end{thebibliography}

\end{document}


